%% sections/06-conclusion.tex

\section{Conclusion}
\label{sec:conclusion}

We have described innovation-cluster-driven rubric amendment as a replicable mechanism
for evidence-based evolution of scoring rubrics applied to creative evaluation. The
mechanism's three stages --- log, cluster, amend --- produce a self-calibrating system
in which the amendment rate declines as the rubric matures and reactivates when novel
creative archetypes introduce quality phenomena the existing anchors cannot capture.

The Marathon corpus demonstrates the mechanism at operational scale: 120+ innovations,
13 rubric versions, 49 sessions, 7 campaigns over 4 days of production use. The key
empirical findings are: (a) all amendment-triggering clusters were thematically
coherent; (b) the innovation rate decayed from 5.7 per session in Campaign 1 to
1.0--2.1 in Campaigns 5--7; (c) Campaign 4's novel siege-resource mechanic reactivated
the amendment rate to 0.57 per session despite the rubric being at version v1.8;
and (d) all 78 universal-scope innovations were adopted into amendments, while
party-specific and adventure-specific innovations served as supporting evidence only.

For practitioners designing evaluation rubrics for generative creative systems, we
recommend: (a) build innovation logging into the evaluation pipeline from session 1;
(b) use a 2+ cluster threshold for amendment proposals; (c) enforce forward-only
amendment application; and (d) expect the mechanism to require ongoing investment
in the dimension that scores the most complex emergent phenomena --- in narrative
games, this is the dimension scoring emotional realization.

The mechanism is available for reuse. The rubric format, innovation log schema, and
amendment proposal template are maintained in the Marathon workshop repository.

The amendment-pipeline methodology bridges two research communities: games research
(FDG, DiGRA) values the playtest and player-behavior findings; the evaluation-methodology
community (NeurIPS D\&B, EMNLP) values the self-calibrating rubric mechanism. We
recommend submission to FDG 2027 as the primary venue, with a parallel shorter-format
position paper targeting NeurIPS D\&B. The FDG framing centers design contribution ---
how the pipeline enables rapid, evidence-grounded iteration on adventure designs and
scoring instruments simultaneously; the NeurIPS D\&B framing centers the evaluation
methodology as a transferable technique applicable to any generative creative system
where session-level quality scoring and corpus growth co-evolve.
